\chapter{Datatypes}
%\appendix{Datatypes }
%
\section {Introduction}
For the purpose of I/O and communication, datatypes can be considered
as abstract representations of the state of an object.  Usually this
is represented by the memory locations where data is stored.  In the
Loci framework, we gave to specify how to save and restore the state
of new object types explicitly to facilitate interprocessor
communication in heterogeneous environments or portable file I/O.
This information is provided to Loci using the traits mechanism.

In Loci, the datatype information is encapsulated in {\tt
  data\_schema\_traits} template class.  A user will need to provide a
specialized template for his/her own datatypes before they can be
used with Loci containers such as {\tt store}.
%
\section {Classification of Datatypes} 

Datatypes could be classified according to the their relationship
between computer memory representation and data which they hold.  The
basic distinction depends on the way C++ stores the object in memory.
If the object is represented in a continuous segment of memory with a
fixed layout, then then the memory layout is consistent with the
schemas used to send messages or write data to files.  If, instead the
objects data is scattered through memory or has a size that is
determined at run time, then the object must be serialized before it
can be sent as a message or placed in file storage.  An object whose
memory layout is compatible with the contiguous memory schema uses an
{\tt IDENTITY\_SCHEMA\_CONVERTER} for serialization.  Other objects
must also specify a schema converter.  The specifics of this will be
specified in the below examples:
%
\begin {verbatim}
typedef struct My_Type_t {
    int       iScalar;
    float     fScalar;
    char      cScalar;
    double    dScalar;
} My_Type1;

typedef struct My_Type_t {
    int        iScalar;
    float      fScalar;:
    char       cScalar;
    double    *dScalar;
} My_Type2;
\end{verbatim}
%
For {\em My\_Type1} C/C++ guarantees that the members of a structure
are stored in contiguous memory locations with some padding, so the
bitwise copy of {\em My\_Type1} will pack the data correctly, whereas
for {\em My\_Type2} the address of {\em dScalar} and not the data will
be copied. In the first case, the memory schema and the file schema is
handled by an {\bf Identity Function} called the {\tt
  IDENTITY\_SCHEMA\_CONVERTER }

In the second case, we need to explicitly convert between memory
schema and the file schema by specifying {\tt USER\_DEFINED\_CONVERTER}.

It should be mentioned here, that pointers are not the only source of
difference between these two datatype. STL container like vector,
list, queues, maps and virtual classes, all would need user to
serialize the state they contains, therefore, they also fall under the
category of user defined converter type.


\section {Predefined Datatypes in Loci}
Loci provides some of the frequently used datatypes so that user need
not write for themselves.These are
%
%
\begin{enumerate}
\item {\bf Atomic Datatype}
As the name implies, these datatypes are indivisible types. These type
are supported by all machines. These are building blocks for all other
datatypes. All native C/C++ datatypes are atomic datatypes
in Loci. The following table provides all the atomic datatypes support
by the Loci.
%
\begin{center}
\begin{tabular} [h] {|l|l|} \hline
Loci Datatype & Native C++ Datatype \\ \hline \hline
BOOL & bool \\ \hline
CHAR & signed char \\\hline
UNSIGNED\_CHAR & unsigned char \\\hline
SHORT & short \\\hline
UNSIGNED\_SHORT & unsigned short \\\hline
INT & int \\\hline
UNSIGNED & unsigned \\\hline
LONG & long \\\hline
UNSIGNED\_LONG & unsigned long \\\hline
FLOAT & float \\\hline
DOUBLE & double \\\hline
\end{tabular}
\end{center}
\item {\bf Array class}
Loci, provides template version of constant size array
\begin{verbatim}
template <class T, unsigned int n>
class Array {
   public:
   .
   .
   private:
      T  x[n];
};
\end{verbatim}
\item {\bf Vector class}
Loci, has 2D/3D vesions of vector class(mathematical vectors)
\begin{verbatim}
 template<T>
 struct vector3d {
      T x,y,z ;
 }

 template<T>
 struct vector2d {
      T x,y ;
 }
\end{verbatim}
\item {\bf STL containers parameterized by types using the {\tt
      IDENTITY\_SCHEMA\_CONVERTER}} 

All standard STL containers with predefined identity schema types are supported.
\begin{enumerate}
\item vector
\item list
\item queue
\item set
\item map
\end{enumerate}
\end{enumerate}
%
\section {The {\tt nameList} Datatype}
%
Many applications need to let users give a list of names in the fact
database, for example a list of variables to include in plot
output, or a list of boundary conditions to apply a setting to.  Loci
provides the {\tt Loci::nameList} datatype for this.  It is defined in
{\tt Tools/nameList.h}, which is included automatically by the
{\tt Loci} header (through {\tt Loci\_types.h}).  A {\tt nameList} is an ordered
list of names, stored as a {\tt std::vector<std::string>}.
%
\subsection {Input Syntax}
%
A {\tt nameList} is read in ASCII form (for example, from a {\tt .vars}
file) as a comma-separated list of names.  The list can be given in
one of two forms:
\begin{itemize}
\item {\bf Bracketed form:} the list is enclosed in square brackets,
  e.g. {\tt [ pg, t, v ]}.  An empty list is written as {\tt []}.
\item {\bf Freeform:} the names are given without brackets,
  e.g. {\tt pg, t, v}.  The list ends at the first name that is not
  followed by a comma.
\end{itemize}
Each entry must be a valid name: it must start with a letter or an
underscore ({\tt \_}), and may contain only letters, digits, and
underscores.  Whitespace (including newlines) between names and
commas is ignored, so a long list can be spread over several lines.
Numbers, quoted strings, and names containing other characters (such
as {\tt -} or {\tt .}) are not accepted.  The following are all valid
entries in a {\tt .vars} file:
\begin{verbatim}
plot_output: P, t, v
plot_output_exclusive: [ pg, t, v ]
boundary_list: [ inflow_1, outflow_1,
                 wall_top, wall_bottom ]
empty_list: []
\end{verbatim}
%
If the input is malformed, for example a missing name after a comma
({\tt [a, , b]}), a missing closing bracket ({\tt [a, b}), or an entry
that is not a valid name, reading the list throws an {\tt exprError}
with a syntax error message.  Note that an empty list can only be
given in bracketed form; an empty freeform entry is a syntax error.
%
\subsection {Output Format}
%
When a {\tt nameList} is written to a stream it is always printed in
bracketed form, with the names separated by a comma and a space,
e.g. {\tt [pg, t, v]}.  An empty list is printed as {\tt []}.
%
\subsection {Using a {\tt nameList} in a Loci Program}
%
To use a {\tt nameList} as a fact, declare its type in the usual way,
usually as a {\tt param}.  It is a good idea to give a default rule
that provides an empty list, so that the user does not have to
specify the variable:
\begin{verbatim}
$type plot_output param<Loci::nameList> ;

$rule default(plot_output) {
  $plot_output.clear() ;
}
\end{verbatim}
%
The class provides a small interface for accessing the list:
\begin{itemize}
\item {\tt begin()} and {\tt end()} return {\tt const} iterators over the
  names (as {\tt std::string} values), so a {\tt nameList} can be used
  in a range-based {\tt for} loop.
\item {\tt empty()} returns {\tt true} if the list has no entries.
\item {\tt clear()} removes all entries from the list.
\end{itemize}
The names are kept in the order the user gave them.  For example, the
following code loops over the names in a {\tt nameList}:
\begin{verbatim}
param<Loci::nameList> plot_output ;
plot_output = facts.get_variable("plot_output") ;

if(!(*plot_output).empty()) {
  for(const auto &name : *plot_output) {
    cout << "plotting variable " << name << endl ;
  }
}
\end{verbatim}
%
\subsection {Serialization}
%
Because a {\tt nameList} holds a variable number of strings, it is not
stored in contiguous memory and uses a {\tt USER\_DEFINED\_CONVERTER}.
Loci serializes it as a {\tt char} buffer through its ASCII stream
representation, using the {\tt StringStreamConverter} class.  This
specialization is already provided in {\tt Loci\_types.h}, so no extra
work is needed to communicate a {\tt nameList} between processors or
write it to a file:
\begin{verbatim}
template<> struct data_schema_traits<nameList> {
  typedef USER_DEFINED_CONVERTER Schema_Converter ;
  typedef char Converter_Base_Type ;
  typedef StringStreamConverter<nameList> Converter_Type ;
} ;
\end{verbatim}
%
\section {The {\tt options\_list} Datatype}
%
Some inputs are too complex to describe with a single number or name.
For example, a boundary condition has a type (such as {\tt outflow})
and may also need its own settings (such as an outflow pressure).  An
initial condition sets several flow variables at once, each with its
own units.  For inputs like these, Loci provides the {\tt
  Loci::options\_list} datatype, defined in {\tt
  Tools/options\_list.h}.  Like {\tt nameList}, it is included
automatically by the {\tt Loci} header.  An {\tt options\_list} is a
set of named options, each of which holds a typed value.  The values
can be numbers with units, names, strings, booleans, lists, or
function-like expressions that hold their own nested options.
%
\subsection {Input Syntax}
%
An {\tt options\_list} is written as a comma-separated list of
assignments of the form \verb|option=value|, enclosed in angle
brackets.  For example, the initial conditions for the {\tt
  quickTest/FVM} cylinder test are given in the {\tt .vars} file as:
\begin{verbatim}
initialConditions  : < T = 300K, p = 1atm, M=0.3>
\end{verbatim}
The rules for the list are:
\begin{itemize}
\item Option names follow the same rules as names in a {\tt nameList}:
  they start with a letter or underscore, and contain only letters,
  digits, and underscores.
\item An option given without a value, e.g. \verb|<prescribed>|, is a
  boolean option with the value {\tt true}.
\item Each option may appear only once in a list.  Assigning an option
  twice, e.g. \verb|<a=1,a=2>|, is an error.
\item Whitespace, newlines, and C++ style \verb|//| comments are
  allowed between entries, so a long list can be spread over several
  lines and commented.
\item An empty list is written as \verb|<>|.
\end{itemize}
%
Each value can be one of the following types.  The type is detected
from the input, and can be queried in the program using the
{\tt option\_value\_type} enumeration names shown in the table:
\begin{center}
\begin{tabular}{|l|l|l|} \hline
Type & Example & Description \\ \hline \hline
{\tt REAL} & \verb|p=1.2e5| & a number without units \\ \hline
{\tt UNIT\_VALUE} & \verb|p=9.e4 Pa|, \verb|T=300K| & a number followed
by units \\ \hline
{\tt NAME} & \verb|limiter=barth| & a name \\ \hline
{\tt STRING} & \verb|file="data/in.dat"| & a quoted string \\ \hline
{\tt BOOLEAN} & \verb|flag=$true|, \verb|flag=$false| & a boolean value \\ \hline
{\tt LIST} & \verb|normal=[1,0,0]| & a list of values in square
brackets \\ \hline
{\tt FUNCTION} & \verb|outflow(p=9.e4 Pa)| & a name followed by a list
of arguments \\ \hline
\end{tabular}
\end{center}
%
A few notes on these value types:
\begin{itemize}
\item Units are given after the number, with or without a space
  (\verb|300K| and \verb|300 K| are the same).  Compound units are built
  with \verb|*|, \verb|/|, and parentheses, e.g. \verb|kg/s/m/m| or
  \verb|one/(Pa*s)|.  The units are checked and converted when the
  program reads the value (see below).
\item The elements of a {\tt LIST} and the arguments of a {\tt FUNCTION}
  can be any of the value types, including nested lists and functions.
  They can also be assignments of the form \verb|name=value|, which
  have type {\tt NAME\_ASSIGN}.  Function arguments are usually written
  as assignments, so that they can be read as an {\tt options\_list}
  of their own (see below).
\item A number can carry derivative information for automatic
  differentiation, written with a caret: \verb|3^2| gives the value 3
  with a first derivative of 2, \verb|3^2^^4| also gives a second
  derivative of 4, and \verb|5^[2,3,4]| gives several first
  derivatives.  This is only needed for sensitivity analysis.
\end{itemize}
%
Any syntax error while reading the list, such as a missing \verb|<|,
a missing comma between options, or an unterminated string, throws
a {\tt StringError} that names the option where the error occurred.
%
\subsection {Output Format}
%
When an {\tt options\_list} is written to a stream it is printed in the
same angle-bracket syntax, e.g. \verb|<M=0.3,T=300 K,p=1 atm>|, so it
can be read back in.  The options are printed in alphabetical order of
their names, not in the order they were given.
%
\subsection {Reading Options in a Program}
%
An {\tt options\_list} is usually declared as a {\tt param} so the user
can give it in the {\tt .vars} file.  For example, the {\tt
  quickTest/FVM} solver declares its initial conditions as:
\begin{verbatim}
$type initialConditions param<Loci::options_list> ;
$type initialConditionRegions param<Loci::options_list> ;
\end{verbatim}
The main member functions for reading the contents of the list are:
\begin{itemize}
\item {\tt optionExists(name)} returns {\tt true} if the option was
  given.
\item {\tt getOptionValueType(name)} returns the type of the option's
  value ({\tt REAL}, {\tt UNIT\_VALUE}, {\tt NAME}, and so on), or
  {\tt NOT\_ASSIGNED} if the option was not given.
\item {\tt getOptionNameList()} returns a {\tt std::list} of the option
  names in the list, in alphabetical order.
\item {\tt getOption(name, value)} copies the option's value into
  {\tt value}.  It is overloaded for {\tt bool}, {\tt double}, {\tt
    std::string} (for names and strings), {\tt UNIT\_type}, and {\tt
    options\_list::arg\_list} (for lists).  For functions, use {\tt
    getOption(name, funcname, args)}, which returns both the function
  name and its argument list.
\item {\tt checkOption(name, value)} returns {\tt true} if the option
  is a name equal to {\tt value}.
\item {\tt getOptionUnits(name, units, value)} reads a number and
  converts it to the given units (see below).
\end{itemize}
%
If {\tt getOption} or {\tt getOptionUnits} is called for an option that
does not exist, it throws a {\tt StringError}.  If the option exists
but has a different type than requested, {\tt getOption} does not
throw an error and does not return a useful value (a warning is printed
only in debug builds).  It is therefore best to check the type first
with {\tt getOptionValueType} when more than one type is allowed.
%
\subsubsection {Units}
%
Most physical inputs should be read with {\tt getOptionUnits}, which
handles both kinds of numbers:
\begin{itemize}
\item If the value has units ({\tt UNIT\_VALUE}), it is converted to the
  requested units.  If the units are not compatible (for example, a
  length given where a pressure is expected) a {\tt StringError} is
  thrown.
\item If the value has no units ({\tt REAL}), it is assumed to already
  be in the requested units and is returned unchanged.
\end{itemize}
For example, the {\tt quickTest/FVM} outflow boundary condition reads
its pressure with:
\begin{verbatim}
$BC_options.getOptionUnits("p","Pa",$p_BC) ;
\end{verbatim}
so the user can give \verb|p=9.e4 Pa|, \verb|p=0.888 atm|, or simply
\verb|p=9.e4|.
%
\subsubsection {Vectors}
%
{\tt getOptionUnits} also has overloads that read a {\tt vector3d}.
A vector can be given in three ways:
\begin{itemize}
\item as a list of three components, e.g. \verb|normal=[1,0,0]| or
  \verb|center=[1 m, 0 m, 2 cm]|;
\item as a single number, which gives the vector $(v,0,0)$, e.g.
  \verb|M=0.3| for a flow in the $x$ direction;
\item in polar form, \verb|polar(r, theta, eta)|, where the angles are
  in degrees unless units are given, e.g.
  \verb|u=polar(100 m/s, 10 deg, 0 deg)|.
\end{itemize}
An optional scale argument multiplies any component that was given
without units.  This is useful for applying a reference length, e.g.
{\tt getOptionUnits("center","m",pt,Lref)}.
%
\subsubsection {Nested Options}
%
The arguments of a {\tt FUNCTION} value, or the elements of a {\tt LIST},
can be turned into an {\tt options\_list} of their own by passing the
argument list to the {\tt Input} member function.  This is how
options with their own settings are read.  For example, the following
code reads an option of the form \verb|default=state(p=1e5, T=300K, u=0)|:
\begin{verbatim}
options_list::arg_list value_list ;
string fname ;
ol.getOption("default",fname,value_list) ;
if(fname != "state") {
  cerr << "expecting function 'state', found '" << fname << "'" << endl ;
  Loci::Abort() ;
}
options_list state ;
state.Input(value_list) ;   // state now holds p, T, and u
\end{verbatim}
%
\subsubsection {Restricting the Allowed Options}
%
By default any option name is accepted.  To catch spelling mistakes, an
{\tt options\_list} can be constructed with a colon-separated list of
the allowed names, e.g. {\tt options\_list("T:p:rho:u:M")}.  When this
list is read from a stream, any other option name causes an error.
%
\subsection {Example: Boundary Conditions}
%
The most common use of {\tt options\_list} is to give boundary
conditions in the {\tt boundary\_conditions} variable.  In this list,
each option name is the name of a boundary surface in the grid file,
and each value is the boundary condition to apply to it.  A boundary
condition with no settings is given as a name, and one with settings is
given as a function whose arguments are its options.  This example is
from the {\tt fixedMass} test in {\tt quickTest/FVM}:
\begin{verbatim}
boundary_conditions: <
 BC_1=reflecting,// two symmetry planes
 BC_2=reflecting,
 BC_3=reflecting,
 BC_4=reflecting,
 BC_5=fixedMass(mdot=144.0 kg/sec, T0=400K),
 BC_6=outflow(p=9.e4 Pa)>
\end{verbatim}
%
The function {\tt setupBoundaryConditions} (in {\tt
  src/System/FVMStuff.cc}) reads this list when the grid is set up and
turns it into facts that rules can use directly:
\begin{itemize}
\item For each boundary condition type {\tt X}, a constraint {\tt X\_BC}
  is created that holds the boundary faces using that condition (for
  example {\tt outflow\_BC} and {\tt fixedMass\_BC}).
\item The settings of each boundary are stored in {\tt BC\_options}, a
  {\tt store<Loci::options\_list>} defined over the boundary surfaces.
  Boundary faces reach it through the {\tt ref} map.
\item For each option name {\tt Y} used by any boundary, a constraint
  {\tt Y\_BCoption} is created that holds the boundary surfaces that set
  that option (for example {\tt p\_BCoption} and {\tt mdot\_BCoption}).
\item Every boundary surface in the grid must appear in {\tt
  boundary\_conditions}, otherwise the setup aborts with an error.
\end{itemize}
%
With these facts, reading a boundary condition setting only needs
a short rule that is constrained to the surfaces that set the option.
For example, the {\tt fixedMass} boundary condition in {\tt
  quickTest/FVM/bcFixedMass.loci} reads its settings with:
\begin{verbatim}
$rule pointwise(mdot_BC<-BC_options),constraint(mdot_BCoption) {
  $BC_options.getOptionUnits("mdot","kg/s",$mdot_BC) ;
}
$rule pointwise(T0_BC<-BC_options),constraint(T0_BCoption) {
  $BC_options.getOptionUnits("T0","kelvin",$T0_BC) ;
}
\end{verbatim}
and the face rules that apply the boundary condition use the
constraint {\tt fixedMass\_BC}, reaching the surface values through
the {\tt ref} map:
\begin{verbatim}
$rule pointwise(T0Ref<-ref->T0_BC) {
  $T0Ref = $ref->$T0_BC ;
}
\end{verbatim}
%
Because a mistake in the boundary conditions may only show up late in
a run, it is a good idea to check the options when the case is set up.
The {\tt quickTest/FVM} solver does this by registering a checker
class for each boundary condition type (see {\tt readGrid.h} and
{\tt bcOutflow.loci}).  The checker lists the options it accepts and
uses the functions above to check the required options and their
units.
%
\subsection {Example: Initial Conditions}
%
A uniform initial condition can be given as a single {\tt options\_list}
that sets the flow state, for example:
\begin{verbatim}
initialConditions   : < p = 9.e4 Pa,  T = 300 K,  M=.4>
\end{verbatim}
The {\tt quickTest/FVM} solver (in {\tt solverIC.loci}) checks which
options are given and reads each one with its units.  Any two of {\tt
  rho}, {\tt p}, and {\tt T} may be given, along with either the
velocity {\tt u} or the Mach number {\tt M}, both read as vectors:
\begin{verbatim}
const Loci::options_list finput = $initialConditions ;
if(finput.optionExists("u"))
  finput.getOptionUnits("u","m/s",u) ;
\end{verbatim}
%
A non-uniform initial condition uses nested options.  The {\tt
  initialConditionRegions} variable defines one or more named states
using a {\tt state} function, and a {\tt regions} list of geometric
tests that select which cells get which state.  This example is from
the shock tube test in {\tt quickTest/FVM}:
\begin{verbatim}
initialConditionRegions   : <
driver = state(p=1e6, T=300K, u=0),
default = state(p=1e5, T=300K, u=0),
regions=[leftPlane(point=[0,0,0],normal=[1,0,0],composition=driver)]
>
\end{verbatim}
Here every cell starts in the {\tt default} state, except for cells
behind the plane through the origin with normal $(1,0,0)$, which get the
{\tt driver} state.  To read this, the solver loops over {\tt
  getOptionNameList()} to find the named states.  It reads each state
function into its own {\tt options\_list} as shown in the nested
options example above.  It then reads the {\tt regions} option as an
{\tt arg\_list}, each element of which is a {\tt FUNCTION} whose name
selects the geometric test ({\tt leftPlane}, {\tt inBox}, {\tt
  inSphere}, {\tt inCylinder}, or {\tt inCone}) and whose arguments
are read with {\tt Input} to get the test's settings.
%
\subsection {Serialization}
%
Like {\tt nameList}, an {\tt options\_list} uses a {\tt
  USER\_DEFINED\_CONVERTER} that serializes it through its ASCII
stream representation.  The specialization is provided in {\tt
  Loci\_types.h}, so an {\tt options\_list} can be stored in a {\tt
  param} or {\tt store} and communicated between processors without any
extra work:
\begin{verbatim}
template <> struct data_schema_traits<options_list> {
  typedef USER_DEFINED_CONVERTER Schema_Converter ;
  typedef char Converter_Base_Type ;
  typedef StringStreamConverter<options_list> Converter_Type ;
} ;
\end{verbatim}
\section {Creating your own compound datatypes}
%
Loci, uses {\tt data\_schema\_traits} template class to determine
the datatype of an object. The generic behavior of template class is
not suitable for identifying or creating new datatype, so we use {\bf
  template Specialization } technique to customize or create new
datatypes . This {\tt data\_schema\_traits} class has one static
member function {\tt get\_type()} in which user specifies the
information about new datatype. A general skeleton for creating new
datatype look as follows
\begin{verbatim}
namespace Loci {
   // Skelton for datatype having identity schema
   template <>
   struct data_schema_traits <My_New_Type1 > {
         typedef IDENTITY_CONVERTER Schema_Converter;
         static DatatypeP get_type() {
                CompoundDatatypeP cmpd = CompoundFactory(My_New_Type1());
                LOCI_INSERT_TYPE(cmpd, My_New_Type1, member);
                   .
                   .
                return DatatypeP(cmpd);
         }
   };
}
\end{verbatim}
%
Where {\bf CompoundFactory} is one of the software design pattern for creating a 
new compound datatype object. Since the allocation of this object is done inside
the functions, the question will always arise, who is responsible for deleting
the object ? To destroy the objects when they are not needed, we use reference counting and the 
{\tt CompoundDatatypeP} class stands for {\em Reference Counting} version of
{\tt CompoundDatatype} class.
%
\par If your C/C++ structures contains members of predefined Loci datatypes, then it
is fairly easy to create corresponding Loci datatype. For example
\begin{verbatim}
 typedef struct My_Compound_Type_t {
        float                       fScalar;      
        vector3d<double>            vect3d;
        Array<double,2>             array1d;
        Array<Array<double,2>,4>    array2d;
 } My_Compound_Type;

 namespace Loci {
  template <> struct data_schema_traits <My_Compound_Type > {
     typedef IDENTITY_CONVERTER Schema_Converter;
     static DatatypeP get_type() {
         CompoundDatatypeP cmpd = CompoundFactory(My_Compound_Type());
         LOCI_INSERT_TYPE(cmpd, My_Compound_Type, fScalar);
         LOCI_INSERT_TYPE(cmpd, My_Compound_Type, vect3d);
         LOCI_INSERT_TYPE(cmpd, My_Compound_Type, array1d);
         LOCI_INSERT_TYPE(cmpd, My_Compound_Type, array2d);
         return DatatypeP(cmpd);
      }
    };
  }
\end{verbatim}
\section{Creating User Defined Datatype}
\par As explained earlier, whenever the memory allocation in any object is
not contiguous, it is the users responsibility to serialize the
data contained in the objects.  The skeleton of user defined schema type
will look as follows
\begin{verbatim}
   // Skelton for datatype having user defined schema
   template <>
   struct data_schema_traits <My_New_Type2 > {
        typedef USER_DEFINED_CONVERTER Schema_Converter;
        typedef char  Converter_Base_Type;
        typedef MyObject_SchemaConverter   Converter_Type;
   };
\end{verbatim}
\par {\em Note : {\tt Converter\_Base\_Type} could be any datatype with identity schema}
\par Following steps must be taken in order to define your own datatype
for Loci
\begin{enumerate} 
\item Specialize the {\tt data\_schema\_traits} class
\begin{itemize}
\item Specialize {\tt data\_schema\_traits} template class with your
class
\item \par declare in {\tt data\_schema\_traits} class
\begin{center}
{\tt typedef USER\_DEFINED\_CONVERTER Schema\_Converter}
\end{center}
\item specify what datatype will be used in conversion. This should be 
datatype with identity schema defined which means, that we can use any 
valid Loci atomic or compound datatype.
\item specify the class which has the responsibility of conversion (
Serialize class )
\end{itemize}
\item Specifying Serialize class
\begin{itemize}
\item Specify object reference in the constructor.
\item {\tt getSize()} member function returns the number of atomic
  datatypes used in this object.(It is not the size of the object in
  bytes)
\item {\tt getState()} member function gets the state of an
object into a contiguous buffer.
\item {\tt setState()} member function sets the state of an
object from a contiguous buffer.
\end{itemize}
\item Overload input/output stream functions.
Both atomic and compound datatypes have already been overloaded with
input/output streams in Loci.{\bf It is required that these function 
are overloaded even if the a user doesn't have intention of
using them.}
\end{enumerate}
%
Now we shall give one simple example to show how things work. We define
one structure with STL list inside it. Since list may not have contiguous
memory, we define it is defined with user defined schema
\begin{verbatim}

  namespace Loci {
  /////////////////////////////////////////////////////////////////////////////////
  // This is an example of conventional C/C++ structure 
  /////////////////////////////////////////////////////////////////////////////////

     struct My_Type {
      list<int>  alist;
      friend ostream& operator << (ostream &, const My_Type &);
      friend istream& operator >> (istream &, My_Type &);
  };

  //------------------------------------------------------------------------------/
  class My_Type_SchemaConverter;    // Forward Declaration of class
  //------------------------------------------------------------------------------/
  // Specialize the data_schema_traits class with "My_Type" class
  //------------------------------------------------------------------------------/

  template <>
  struct data_schema_traits<My_Type> {
    // This class has user defined schema
    typedef USER_DEFINED_CONVERTER Schema_Converter ;

    // Since list contains "int" we use it directly for our converion
    typedef int Converter_Base_Type ;
   
    // Here we specify the class used for serialization/deserialization 
    // purpose
    typedef My_Type_SchemaConverter Converter_Type ;
  };

  //------------------------------------------------------------------------------/
  // Define a class which has the responsibity of serialization and deserialization 
  // of "My_Type" class
  //------------------------------------------------------------------------------/

  class My_Type_SchemaConverter {
    // For the schema converter, we always store a reference to the object
    // we are converting schmata for.
    My_Type &RefObj ;
    public:
        explicit My_Type_SchemaConverter(My_Type &new_obj): RefObj(new_obj) {}
        //
        // This member function returns number of elements of type defined
        // in Converter_Base_Type. It is not the size in bytes.
        //
        int getSize() const {
            return RefObj.alist.size() ;
        }
      
        // Get the state of an object "RefObj" into an array and also size of
        // array. This is a serialization step. 
        void getState(int *buf, int &size) {
             size = getSize() ;
             int ii=0;

             list<int> :: const_iterator ci;
             list<int> :: const_iterator begin = RefObj.alist.begin();
             list<int> :: const_iterator end   = RefObj.alist.end();
             for(ci = begin; ci != end; ++ci) 
                 buf[ii++] = *ci;
        }
        //
        // From a given array, construct the object. This is "Deserialization Step"
        //
        void setState(int *buf, int size) {
             RefObj.alist.clear();
             list<int> :: iterator ci;
             for(int i=0;i<size;++i)
                 RefObj.alist.push_back(buf[i]);
        }
  };
  }
  //------------------------------------------------------------------------------/
\end{verbatim}
%
%
\section{Inner Details about Compound Datatype}
Compound datatypes are similar to structures in C/C++. These datatypes
are a collection of heterogeneous atomic or fixed sized array datatypes. 
Every member of these datatype has a unique name within the datatype and 
they occupy  non-overlapping memory locations.
%
%
\par In Loci, this datatype is declared in {\tt CompoundType}
class. The corresponding counted pointer class is {\tt
CompoundDatatypeP}.
%
\par A new member can be inserted into the new datatype in either way
\begin{itemize}
\item Using member function of CompoundType class
\begin{center}
insert( member\_name, offsetof(type, member-designator), member\_datatype);
\end{center}
where {\tt offsetof} is a standard C/C++ function which provides offset of any
member (designated by member-designator) in C/C++ structure (designated by type).
{\em member\_datatype} could be any valid Loci datatype.
%
\item Using predefined macro
\begin{center}
LOCI\_INSERT\_TYPE( compound\_object, compound\_class, insert\_member);
\end{center}
\par where {\em compound\_object} is the compound datatype for {\em compound\_class}
and {\em insert\_member} is the required member of {\em compound\_class} which
is inserted into new datatype. 
\par In order to use the macro {\em insert\_member} must be a first class object
and and its own type should be identified by {\tt data\_schema\_traits} class.
\end{itemize}
\par In the following sections we shall gives some examples of
creating different Loci datatypes.
%
\subsection {Creating compound datatype with only atomic datatypes}
The following is a very simple C/C++ structure, which contains only native datatypes.
For this structure, we would like to create Loci datatype, which is also described
below.
\begin{verbatim}
 typedef struct My_Compound_Type_t {
        int        iScalar;
        float      fScalar;
        char       cScalar;
        double     dScalar;
 } My_Compound_Type;
\end{verbatim}
\begin{verbatim}
  // data_scheme_traits should be defined in Loci namespace
 namespace Loci {
   // Specialize the class with the new class
   template <> 
   struct data_schema_traits <My_Compound_Type > {

        // Specify that ideneity Schema will be used for this class
        typedef IDENTITY_CONVERTER Schema_Converter;

        // define the member function
        static DatatypeP get_type() {
              // Create a new product "cmpd" from the factory pattern 
              CompoundDatatypeP cmpd = CompoundFactory(My_Compound_Type());
              // Insert a new member into the new compound datatype
              LOCI_INSERT_TYPE(cmpd, My_Compound_Type, iScalar);
              LOCI_INSERT_TYPE(cmpd, My_Compound_Type, fScalar);
              LOCI_INSERT_TYPE(cmpd, My_Compound_Type, cScalar);
              LOCI_INSERT_TYPE(cmpd, My_Compound_Type, dScalar);
             // return pointer to the base class
             return DatatypeP(cmpd);
         }
      };
  }
\end{verbatim}
\subsection{Creating compound datatype with arrays}
In the following example, we have inserted a two dimensional array
into the structure and define corresponding Loci datatype.
\begin{verbatim}
 namespace Loci {
     typedef struct My_Compound_Type_t {
         int     iScalar;
         float   fScalar;
         double  dScalar;
         Array<double,10>          dArray1D;
         Array<Array<double,3>,5>  dArray2D;
     } My_Compound_Type;
    
     template<>
     struct data_schema_traits<My_Compound_Type> {
          typedef IDENTITY_CONVERTER Schema_Converter ;
          static DatatypeP get_type() {
            CompoundDatatypeP ct = CompoundFactory(My_Compound_Type()) ;
            LOCI_INSERT_TYPE(ct,My_Compound_Type, iScalar) ;
            LOCI_INSERT_TYPE(ct,My_Compound_Type, fScalar) ;
            LOCI_INSERT_TYPE(ct,My_Compound_Type, dScalar) ;
            LOCI_INSERT_TYPE(ct,My_Compound_Type, dArray1D) ;
            LOCI_INSERT_TYPE(ct,My_Compound_Type, dArray2D) ;
            return DatatypeP(ct) ;
          }
      } ;
 }
\end{verbatim}
\subsection{Creating compound datatype with nested compound
datatypes}
In this example, we would demonstrate that any member with valid
datatype could be inserted into compound datatype in an hierarchal fashion. There
is no restriction on number of levels used to define a datatype.
\begin{verbatim}
namespace Loci {
   struct Velocity {
     Array<double,3>  comp;
   };
   // Define traits for "Velocity" structure
   template <>
   struct data_schema_traits<Velocity> {
     typedef IDENTITY_CONVERTER Schema_Converter ;
     static DatatypeP get_type() {
       Velocity v;
       return getLociType(v.comp);
     }
   };
   // A structure contains another structure
   struct CellAttrib {
     int      local_id;
     double   density;
     double   pressure;
     Velocity vel;
   };
 
   // Define traits for "CellAttrib" class. Notice that "vel" is a structure
   // and since its type is already defined, we can insert it similar to other
   // members.
   template <>
   struct data_schema_traits<CellAttrib> {
     typedef IDENTITY_CONVERTER Schema_Converter ;
     static DatatypeP get_type() {
       CompoundDatatypeP  cmpd = CompoundFactory( CellAttrib() );
       LOCI_INSERT_TYPE(cmpd, CellAttrib, local_id );
       LOCI_INSERT_TYPE(cmpd, CellAttrib, density  );
       LOCI_INSERT_TYPE(cmpd, CellAttrib, pressure );
       LOCI_INSERT_TYPE(cmpd, CellAttrib, vel );
       return DatatypeP(cmpd);
     }
   };
 }
\end{verbatim}

\section{Array Datatype}
Array datatype consists of homogeneous collection of both compound and 
atomic datatypes. We can defined array datatypes for standard C/C++
arrays. In order to create array datatype, a user need to provide
\begin{itemize}
\item the rank of the array, i.e. number of dimensions 
\par {\em Note At present, Loci can support arrays with maximum rank of
4. This limitation comes from HDF5 library. If the user wants higher ranked
arrays, using {\bf Array class} is one solution}
\item the size of each dimension
\item the datatype of each element of the array
\end{itemize}
\par In Loci, this datatype is declared in {\tt ArrayType} class. The 
corresponding counted pointer class is {\tt ArrayDatatypeP}.
%
\begin{itemize}
\item Create 1 D dimensional array datatype
\begin{verbatim}
dims[0]= 100;
ArrayType  atype(Loci::DOUBLE, 1, dims);
\end{verbatim}
%
\item Create 2 D dimensional array datatype
\begin{verbatim}
dims[0] = 10;
dims[1] = 20;
ArrayType  atype(Loci::DOUBLE, 2, dims);
\end{verbatim}
%
\item Create 3 D dimensional array datatype
\begin{verbatim}
dims[0] = 10;
dims[1] = 20;
dims[2] = 30;
ArrayType  atype(Loci::DOUBLE, 3, dims);
\end{verbatim}
%
\item Create 4 D dimensional array datatype
\begin{verbatim}
dims[0] = 10;
dims[1] = 20;
dims[2] = 30;
dims[3] = 40;
ArrayType  atype(Loci::DOUBLE, 4, dims);
\end{verbatim}
\end{itemize}
%
%
\par For example
\begin{verbatim}
typedef struct My_Compound_Type_t {
    int       iScalar;
    float     fScalar;
    char      cScalar;
    double    dScalar[10][5][2];
}My_Compound_Type;
\end{verbatim}
%
This datatype has multidimensional array, which is not first class
objects. We can use {\tt ArrayDatatype} to specify the Loci Datatype as
\begin{verbatim}
int     rank  = 3;
int     dim[] = {10, 5, 2};
int     sz    = 100*sizeof(double);

My_Compound_Type     type;
CompoundDatatypeP cmpd    = CompoundFactory(My_Compound_Type());
DatatypeP         atom    = getLociType(type.dScalar[0][0][0]);
ArrayDatatypeP    array_t = ArrayFactory(atom, sz, rank, dim);
cmpd->insert("dScalar", offsetof(My_Compound_Type, dScalar), DatatypeP(array_t));
\end{verbatim}

This is definitely cumbersome, Instead of using array in this way, if
we had used
\begin{verbatim}
typedef  Array < double, 10 > Array1D;
typedef  Array < Array1D, 5 > Array2D;
typedef  Array < Array2D, 2 > Array3D;
typedef struct My_Compound_Type_t {
   int             iScalar;
   float           fScalar;
   char            cScalar;
   Array3D         dScalar;
} My_Compound_Type;
\end{verbatim}
then we can use {\tt LOCI\_INSERT\_TYPE} macro to insert {\tt
dScalar} into new compound datatype.

%
